% Part:sets-functions-relations
% Chapter: sets
% Section: reduction

\documentclass[../../../include/open-logic-section]{subfiles}

\begin{document}

\olfileid[fr]{sfr}{siz}{red}

\Needspace{8\baselineskip}\olsection{Réduction}

\begin{editorial}
  Cette section établit la non-dénombrabilité par réduction, en reprenant
  les résultats de la \olref[nen]{sec}. Une autre version, un peu plus
  condensée et correspondant aux résultats de la \olref[nen-alt]{sec},
  figure à la \olref[red-alt]{sec}.
\end{editorial}

Nous avons montré par diagonalisation que $\Pow{\PosInt}$ est
!!{nonenumerable}. Nous disposions déjà d'une démonstration de la
non-dénombrabilité de $\Bin^\omega$, l'ensemble de toutes les suites
infinies de $0$ et de $1$. Voici une autre manière de montrer que
$\Pow{\PosInt}$ est !!{nonenumerable} : établir que \emph{si
$\Pow{\PosInt}$ est !!{enumerable}, alors $\Bin^\omega$ l'est aussi}.
Puisque nous savons que $\Bin^\omega$ n'est pas !!{enumerable},
$\Pow{\PosInt}$ ne peut pas l'être non plus. C'est ce qu'on appelle
\emph{réduire} un problème à un autre : ici, nous réduisons le problème
de l'énumération de $\Bin^\omega$ à celui de l'énumération de
$\Pow{\PosInt}$. Une solution du second problème, à savoir une
énumération de $\Pow{\PosInt}$, fournirait une solution du premier,
c'est-à-dire une énumération de $\Bin^\omega$.

Comment réduire le problème de l'énumération d'un ensemble~$B$ à celui
de l'énumération d'un ensemble~$A$ ? Il faut donner un moyen de
transformer une énumération de~$A$ en une énumération de~$B$. Le plus
simple est de définir une fonction !!{surjective} $f\colon A \to B$.
Si $x_1$, $x_2$, \dots{} énumèrent~$A$, alors $f(x_1)$, $f(x_2)$,
\dots{} énumèrent~$B$. Dans notre cas, nous cherchons donc une fonction
surjective $f\colon \Pow{\PosInt} \to \Bin^\omega$.

\begin{prob}
Montrer que, s'il existe une fonction !!{injective} $g\colon B \to A$
et si $B$ est !!{nonenumerable}, alors $A$ l'est aussi. Pour cela,
expliquer comment utiliser~$g$ pour transformer une énumération de~$A$
en une énumération de~$B$.
\end{prob}

\begin{proof}[Démonstration du {\olref[nen]{thm:nonenum-pownat}} par réduction]
Supposons que $\Pow{\PosInt}$ soit !!{enumerable}. Il en existe alors
une énumération $Z_{1}$, $Z_{2}$, $Z_{3}$, \dots

Définissons $f \colon \Pow{\PosInt} \to \Bin^\omega$ en associant
à chaque $Z$ la suite $s$ telle que $s(n) = 1$ si et seulement si
$n \in Z$, et $s(n) = 0$ sinon.\footnote{Note de traduction :
l'original nomme ici la suite $s_k$ sans définir l'indice $k$.
Nous écrivons $s$ pour la suite associée au seul ensemble $Z$.}
Cela définit bien une fonction : pour tout $Z \subseteq \PosInt$,
chaque $n \in \PosInt$ appartient ou n'appartient pas à~$Z$.
Par exemple, l'ensemble $2\PosInt = \{2, 4, 6, \dots\}$ des entiers
pairs strictement positifs a pour image la suite $010101\dots$ ;
l'ensemble vide a pour image $0000\dots$, et l'ensemble $\PosInt$
lui-même a pour image $1111\dots$.

Cette fonction est aussi !!{surjective} : à toute suite de $0$ et de
$1$ correspond un ensemble d'entiers strictement positifs, celui des
indices auxquels la suite vaut~$1$. Plus précisément, soit
$s \in \Bin^\omega$. Définissons $Z \subseteq \PosInt$ par :
\[
Z = \Setabs{n \in \PosInt}{s(n) = 1}
\]
On a alors $f(Z) = s$, comme on le vérifie à partir de la définition
de~$f$.

Considérons maintenant la liste
\[
f(Z_1), f(Z_2), f(Z_3), \dots
\]
Puisque $f$ est !!{surjective}, chaque élément de $\Bin^\omega$ est
la valeur de~$f$ en un certain argument et doit donc figurer dans
cette liste. Celle-ci énumère ainsi tout~$\Bin^\omega$.

Si $\Pow{\PosInt}$ était !!{enumerable}, $\Bin^\omega$ le serait
donc aussi. Or $\Bin^\omega$ est !!{nonenumerable}
(\olref[nen]{thm:nonenum-bin-omega}). Par conséquent,
$\Pow{\PosInt}$ est !!{nonenumerable}.
\end{proof}

\begin{explain}
Il est facile de se tromper sur le sens de la réduction. Par exemple,
une fonction !!{surjective} $g \colon \Bin^\omega \to B$ ne permet
\emph{pas} d'établir que $B$ est !!{nonenumerable}. Considérons
$g \colon \Bin^\omega \to \Bin$, définie par $g(s) = s(1)$ : elle
associe à une suite de $0$ et de $1$ son premier terme. Elle est
!!{surjective}, car certaines suites commencent par $0$ et d'autres
par $1$. Pourtant, $\Bin$ est fini. Remarquons aussi que la
fonction~$f$ doit être !!{surjective} : sans cette hypothèse,
l'argument ne fonctionne plus, car rien ne garantit que $f(x_1)$,
$f(x_2)$, \dots{} contiennent tous les éléments de~$B$. Par exemple,
\[
h(n) = \underbrace{000\dots0}_{\text{$n$ zéros}}111\dots
\]
définit une fonction $h\colon \PosInt \to \Bin^\omega$, alors que
$\PosInt$ est !!{enumerable}.\footnote{Note de traduction : l'original
affiche seulement un mot fini de $n$ zéros, qui n'appartient pas à
$\Bin^\omega$. Nous le complétons ici par une infinité de $1$.
Cette fonction n'est pas surjective : la suite constante égale à $0$,
par exemple, n'est l'image d'aucun entier.}
\end{explain}

\begin{prob}
Montrer par réduction que l'ensemble de tous les \emph{ensembles de}
couples d'entiers strictement positifs est !!{nonenumerable}.
\end{prob}

\begin{prob}\label{sfr:siz:red:prob:nat-nat}
  Montrer par réduction que l'ensemble~$X$ de toutes les fonctions
  $f\colon \Nat \to \Nat$ est !!{nonenumerable}. Indication : donner
  une fonction surjective de $X$ dans~$\Bin^\omega$.
\end{prob}

\begin{prob}
Montrer par réduction que $\Nat^\omega$, l'ensemble des suites
infinies d'entiers naturels, est !!{nonenumerable}.
\end{prob}

\begin{prob}
Soit $P$ l'ensemble des fonctions de l'ensemble des entiers
strictement positifs dans $\{0\}$, et soit $Q$ l'ensemble des
fonctions \emph{partielles} entre ces mêmes ensembles. Montrer que
$P$ est !!{enumerable} et que $Q$ ne l'est pas. Indication : réduire
le problème de l'énumération de $\Bin^\omega$ à celui de
l'énumération de~$Q$.
\end{prob}

\begin{prob}
Soit $S$ l'ensemble de toutes les fonctions !!{surjective}s de
l'ensemble des entiers strictement positifs dans l'ensemble $\{0,1\}$ :
autrement dit, $S$ contient exactement les fonctions !!{surjective}s
$f\colon \PosInt \to \Bin$. Montrer que $S$ est !!{nonenumerable}.
\end{prob}

\begin{prob}
Montrer que l'ensemble~$\Real$ des nombres réels est !!{nonenumerable}.
\end{prob}

\end{document}
